%! Function Fundamentals — Unit 1, Lesson 1.1 — Guided Notes
% THE in-class document, in the MAIN IDEAS / NOTES shape (LESSON_SHAPE.md §1): the vocabulary
% box, the hook, then ONE two-column guidednotes table. Each row is one idea — a label on the
% left; on the right one or two COMPLETE printed sentences, ONE large pre-drawn display, then a
% two-across grid of numbered problems with real writing room. Problems run 1..18.
% DENSITY: a \blank{} only where a word IS the answer (the direction table); no mid-sentence
% blanks; two problems across; 2–2.6cm of answer space each.
%   I DO   : the teacher reads the definition, marks up the display, works the first problem
%            of each grid (1, 3, 7, 11).
%   WE DO  : the class works the rest of each grid, students holding the pen, every one
%            cold-called.
%   YOU DO : the last row — problems 15–18, alone, for 13 minutes while the teacher circulates.
% The crux is row 4, TWO DIRECTIONS: problems 12 and 13 ask "find f(3)" and "solve f(x) = 3"
% off the same graph and get DIFFERENT answers (3, then 1 and 3). Problem 16 transfers it.
%
% PAGE LOCKSTEP with notes_key/: the vocabulary write lines -> \ansline, \blank -> \ans, and
% the answer argument of every \pcell, \writespace and \labelbox filled. Nothing else differs.
% Inside a cell `\\` ENDS THE ROW — break lines with \par. A row cannot break across pages:
% the figure and EACH grid row are separate sub-rows (`\\` then `& ...`, probgrid*).
\documentclass[10pt]{article}
\usepackage{precalculus-article}
\usepackage{precalculus-boxes}
\newcommand{\TallMath}[1]{$\displaystyle #1\rule[-1.4em]{0pt}{3.2em}$}

\begin{document}

\pageheader{Unit 1, Lesson 1.1}{Guided Notes}

\vspace{0.12in}

\boxguard
\begin{vocabbox}
Fill in each term as we define it together.
\par\vspace{2pt}
\noindent\textbf{\textcolor{plum}{Function notation:}}\\[1pt]
\writeline\\[3pt]
\noindent\textbf{\textcolor{plum}{Evaluate:}}\\[1pt]
\writeline\\[3pt]
\noindent\textbf{\textcolor{plum}{Solve:}}\\[1pt]
\writeline\\[3pt]
\noindent\textbf{\textcolor{plum}{Domain:}}\\[1pt]
\writeline\\[3pt]
\noindent\textbf{\textcolor{plum}{Range:}}\\[1pt]
\writeline\\[3pt]
\noindent\textbf{\textcolor{plum}{Undefined:}}\\[1pt]
\writeline
\end{vocabbox}

\vspace{0.1in}

\boxguard
\begin{hookbox}
The vending machine from yesterday finally gets a name. Call it $P$, and let
$P(b)$ be the price in dollars of the item at button $b$. On the board:
\[ P(\text{B4}) = 1.25 \]

Say that equation out loud as an English sentence. What is the input here ---
and does an input have to be a number?
\writelines{2}
\end{hookbox}

\small
\begin{guidednotes}
% ── ROW 1: function notation — naming the machine ────────────────────────────
\mainidea[Naming the machine]{Notation} &
A function has three parts, and the symbol $f(x)$ names all three at once: $f$ is the
\textbf{name} of the machine, $x$ is the \textbf{input} that goes in the slot, and $f(x)$ is
the \textbf{output} that comes out. It is read ``$f$ of $x$'' --- \textbf{never} ``$f$ times
$x$.'' The parentheses hold the input; they do not multiply.

\vspace{4pt}
A rideshare app charges a flat fee plus a rate per mile: $F(m) = 2.50 + 1.75m$.
\begin{center}
\begin{tikzpicture}[scale=0.95, >=stealth]
  \node[font=\small] (in) at (0,0) {input $m$};
  \node[draw=plum, very thick, fill=lilac, rounded corners=3pt, minimum width=4.2cm,
        minimum height=1.5cm, align=center, font=\small] (box) at (4.2,0)
        {\textbf{\Large\color{plum} $F$}\\ multiply by 1.75, then add 2.50};
  \node[font=\small] (out) at (8.6,0) {output $F(m)$};
  \draw[->, very thick, plum] (in) -- (box);
  \draw[->, very thick, plum] (box) -- (out);
  \node[font=\scriptsize\itshape, above] at (1.55,0.12) {goes in};
  \node[font=\scriptsize\itshape, above] at (7.05,0.12) {comes out};
\end{tikzpicture}
\end{center}
\vspace{2pt}
$m$ stands for \labelbox{2.8cm}{}\qquad $F(m)$ stands for \labelbox{3.2cm}{}
\\
& \notesprompt{Read each statement as a sentence: name the input, the output, and the point on the graph.}
\begin{probgrid}{|Y|Y|}
\pcell{1}{$f(3) = 7$}{2.0cm}{} &
\pcell{2}{$F(4)$, for the rideshare fare above. What is it asking for? Find it.}{2.0cm}{} \\ \hline
\end{probgrid}
\\ \hline
% ── ROW 2: evaluating from a formula — the empty slot ────────────────────────
\mainidea[From a formula]{Evaluating} &
To \textbf{evaluate} a function, drop the input into the slot --- \emph{always inside
parentheses} --- and simplify. Evaluating never changes the rule; it just runs the machine once.
\textbf{Write the parentheses first, empty, then fill them.} That habit is what saves you when
the input is negative.

\vspace{4pt}
\begin{center}
\begin{tikzpicture}[scale=0.95, >=stealth]
  \node[font=\large] (rule) at (0,0) {$f(x) = 5x - 8$};
  \node[font=\large] (slot) at (5.6,0) {$f(\square) = 5(\square) - 8$};
  \draw[->, very thick, plum] (rule) -- node[above, font=\scriptsize\itshape] {empty the slot} (slot);
  \node[font=\large] (fill) at (5.6,-1.5) {$f(-2) = 5(-2) - 8$};
  \draw[->, very thick, plum] (slot) -- node[right, font=\scriptsize\itshape] {then fill it} (fill);
  \node[draw=goldacc, thick, fill=hookbg, rounded corners=2pt, font=\small, align=left]
        at (0,-1.5) {\textbf{The trap:} $(-4)^2$ and $-4^2$\\ are \emph{not} the same number.};
\end{tikzpicture}
\end{center}
\\
& \notesprompt{Evaluate. Show the slot filled before you simplify.}
\begin{probgrid}{|Y|Y|}
\pcell{3}{$f(x) = 5x - 8$. Find $f(3)$.}{2.0cm}{} &
\pcell{4}{$f(x) = 5x - 8$. Find $f(-2)$.}{2.0cm}{} \\ \hline
\end{probgrid}
\\
& \begin{probgrid}*{|Y|Y|}
\pcell{5}{$g(x) = x^2 + 3x$. Find $g(-4)$.}{2.0cm}{} &
\pcell{6}{Compute $(-4)^2$ and $-4^2$. Which one did problem 5 need, and what made the difference?}{2.0cm}{} \\ \hline
\end{probgrid}
\\ \hline
% ── ROW 3: domain and range ──────────────────────────────────────────────────
\mainidea[Every input,]{Domain \& Range} &
The \textbf{domain} is the set of all inputs a function uses; the \textbf{range} is the set
of all outputs it produces, with a repeated output listed once. From a graph, read the domain
left to right and the range bottom to top; write either as an inequality, $0 \le x \le 8$, or
with brackets, $[\,0, 8\,]$ --- a square bracket means the endpoint \textbf{is} included.
\\
& \begin{center}
\begin{tikzpicture}[x=0.55cm, y=0.4cm, >=stealth]
  \draw[linegray, very thin, step=1] (0,0) grid (8,7);
  \draw[->, thick] (0,0) -- (8.6,0) node[below right] {\small $x$};
  \draw[->, thick] (0,0) -- (0,7.6) node[above] {\small $f(x)$};
  \foreach \x in {1,2,3,4,5,6,7,8} \node[below] at (\x,0) {\small \x};
  \foreach \y in {1,3,5,7} \node[left] at (0,\y) {\small \y};
  \draw[very thick, plum] (0,5) -- (2,1) -- (5,7) -- (8,4);
  \fill[plum] (0,5) circle (2.5pt);
  \fill[plum] (8,4) circle (2.5pt);
  % domain bracket under the x-axis
  \draw[very thick, goldacc, |-|] (0,-1.1) -- (8,-1.1);
  \node[below, font=\scriptsize\bfseries\color{goldacc}] at (4,-1.2) {left to right: every $x$ the graph uses};
  % range bracket left of the y-axis
  \draw[very thick, redacc, |-|] (-1.3,1) -- (-1.3,7);
  \node[rotate=90, above, font=\scriptsize\bfseries\color{redacc}] at (-1.5,4) {bottom to top: every height};
\end{tikzpicture}
\end{center}
\vspace{2pt}
Gold bracket: \labelbox{2.0cm}{}\qquad Red bracket: \labelbox{2.0cm}{}
\\
& \notesprompt{State the domain and the range.}
\begin{probgrid}{|Y|Y|}
\pcell{7}{$f = \{(-2, 6),\ (0, 4),\ (3, 6),\ (5, 1)\}$. Domain and range, as lists.}{2.0cm}{} &
\pcell{8}{The graph above: domain and range as inequalities, then again with brackets.}{2.0cm}{} \\ \hline
\end{probgrid}
\\
& \begin{probgrid}*{|Y|Y|}
\pcell{9}{$k(x) = \dfrac{8}{x + 3}$. Which input is forbidden, and why? Write the domain in words.}{2.2cm}{} &
\pcell{10}{In problem 7 the output 6 appears twice. How many times is it listed in the range? Is $f$ still a function?}{2.2cm}{} \\ \hline
\end{probgrid}
\\ \hline
% ── ROW 4: THE CRUX — evaluate vs. solve, the two directions ─────────────────
\mainidea[Evaluate vs.\ solve]{Two Directions} &
Every question in this lesson runs in one of two directions, and the same graph answers both.
\textbf{Read the arrow before you compute.}

\vspace{3pt}
\renewcommand{\arraystretch}{1.35}
\begin{tabular}{|p{0.25\linewidth}|p{0.31\linewidth}|p{0.31\linewidth}|}
\hline
 & \textbf{Find $f(2)$} & \textbf{Solve $f(x) = 2$} \\
\hline
You are given & the \blank{1.8cm} & the \blank{1.8cm} \\
\hline
You are asked for & the \blank{1.8cm} & the \blank{1.8cm} \\
\hline
How many answers & exactly one & \blank{3.0cm} \\
\hline
\end{tabular}
\\
& \begin{center}
\begin{tikzpicture}[x=0.55cm, y=0.4cm, >=stealth]
  \draw[linegray, very thin, step=1] (0,0) grid (8,7);
  \draw[->, thick] (0,0) -- (8.6,0) node[below right] {\small $x$};
  \draw[->, thick] (0,0) -- (0,7.6) node[above] {\small $f(x)$};
  \foreach \x in {1,2,3,4,5,6,7,8} \node[below] at (\x,0) {\small \x};
  \foreach \y in {1,3,5,7} \node[left] at (0,\y) {\small \y};
  \draw[very thick, plum] (0,5) -- (2,1) -- (5,7) -- (8,4);
  \fill[plum] (0,5) circle (2.5pt);
  \fill[plum] (8,4) circle (2.5pt);
  % evaluate path: input 2, walk up, read across
  \draw[->, very thick, goldacc, dashed] (2,-0.05) -- (2,1) -- (0.1,1);
  \node[font=\scriptsize\bfseries\color{goldacc}, below] at (2,-0.95) {start here};
  % solve path: output 3, walk across, drop down twice
  \draw[->, very thick, redacc, dashed] (0.05,3) -- (1,3) -- (1,0.1);
  \draw[->, very thick, redacc, dashed] (1,3) -- (3,3) -- (3,0.1);
  \node[font=\scriptsize\bfseries\color{redacc}, above left] at (-0.2,3.4) {start here};
\end{tikzpicture}
\end{center}
\vspace{2pt}
\textbf{Gold path} (input 2, up, across): $f(2) = $ \labelbox{1.2cm}{}\quad
This is \labelbox{2.3cm}{}\par\vspace{3pt}
\textbf{Red path} (output 3, across, down): $x = $ \labelbox{1.6cm}{}\quad
This is \labelbox{2.3cm}{}
\\
& \fcolorbox{plum}{lilac}{\parbox{\dimexpr\linewidth-2\fboxsep-2\fboxrule\relax}{\centering
\textbf{Evaluate hands you the input and asks for the output --- exactly one answer.
Solve hands you the output and asks for the input(s) --- possibly several, and several is
still a function.}}}

\vspace{3pt}
\textbf{In your own words:} how do you tell an \emph{evaluate} question from a \emph{solve}
question before you compute anything?
\writespace{1.4cm}{}
\\
& \notesprompt{The same graph, both directions. Say which direction each one runs.}
\begin{probgrid}{|Y|Y|}
\pcell{11}{Find $f(6)$.}{2.0cm}{} &
\pcell{12}{Find $f(3)$.}{2.0cm}{} \\ \hline
\end{probgrid}
\\
& \begin{probgrid}*{|Y|Y|}
\pcell{13}{Solve $f(x) = 3$. Compare with problem 12: is it the same question? Is it the same answer?}{2.0cm}{} &
\pcell{14}{Solve $f(x) = 5$. Three inputs share one output --- is $f$ still a function?}{2.0cm}{} \\ \hline
\end{probgrid}
\\ \hline
% ── LAST ROW: YOU DO — alone, while the teacher circulates (13 min) ──────────
\mainidea[You do, alone]{Table \& Formula} &
Two functions given only by a table, $s$ and $p$, and two given by a formula. The arrows
show $p$ running forward: each input on the left goes to its output on the right.
\\
& \begin{center}
\begin{tikzpicture}[scale=1.0, >=stealth]
  \node[font=\normalsize, anchor=east] at (2.4,1.0) {%
    \renewcommand{\arraystretch}{1.3}%
    \begin{tabular}{|c|c|c|c|c|c|}
    \hline $x$ & 1 & 2 & 3 & 4 & 5 \\ \hline $s(x)$ & 12 & 9 & 9 & 6 & 3 \\ \hline
    \end{tabular}};
  \node[font=\normalsize, anchor=east] at (2.4,-1.0) {%
    \renewcommand{\arraystretch}{1.3}%
    \begin{tabular}{|c|c|c|c|c|}
    \hline $x$ & 1 & 2 & 3 & 4 \\ \hline $p(x)$ & 9 & 6 & 6 & 3 \\ \hline
    \end{tabular}};
  \begin{scope}[xshift=4.6cm]
    \node[font=\small\bfseries\color{plum}] at (0,2.0) {$p$: inputs};
    \node[font=\small\bfseries\color{plum}] at (3,2.0) {outputs};
    \foreach \x/\y in {1/1.2, 2/0.4, 3/-0.4, 4/-1.2} \node[font=\small] (i\x) at (0,\y) {\x};
    \foreach \v/\y/\n in {9/1.2/a, 6/0/b, 3/-1.2/c} \node[font=\small] (o\n) at (3,\y) {\v};
    \foreach \a/\b in {i1/oa, i2/ob, i3/ob, i4/oc} \draw[->, thick, plum] (\a) -- (\b);
  \end{scope}
\end{tikzpicture}
\end{center}
\\
& \notesprompt{On your own. Show your work.}
\begin{probgrid}{|Y|Y|}
\pcell{15}{$f(x) = 3x - 7$. Find $f(5)$ and $f(-1)$. Fill the slot first.}{2.2cm}{} &
\pcell{16}{From the $s$ table: find $s(4)$, then solve $s(x) = 9$. Which one was the \emph{evaluate} question?}{2.2cm}{} \\ \hline
\end{probgrid}
\\
& \begin{probgrid}*{|Y|Y|}
\pcell{17}{State the domain and range of $p$. Then run $p$ \emph{backward}: feed it the output 6 and ask for the input. Is that backward rule a function?}{2.4cm}{} &
\pcell{18}{$m(x) = \dfrac{20}{x - 6}$. Which input is forbidden, and why? Write the domain of $m$ in words.}{2.4cm}{} \\ \hline
\end{probgrid}
\\ \hline
\end{guidednotes}

% ── THE NOTES END AT THE YOU DO ROW; AP Practice and the Homework follow ──────

\end{document}
